\documentclass[border=10pt]{standalone}
\usepackage[no-math]{fontspec}
\usepackage[brazil]{babel}
\usepackage{amsmath,amssymb}
\IfFileExists{C:/Windows/Fonts/Lato_Regular.ttf}{%
  \setmainfont{Lato}[
    Path = C:/Windows/Fonts/,
    UprightFont = Lato_Regular.ttf,
    BoldFont = Lato_Bold.ttf,
    ItalicFont = Lato_Italic.ttf,
    BoldItalicFont = Lato_BoldItalic.ttf
  ]
  \setsansfont{Lato}[
    Path = C:/Windows/Fonts/,
    UprightFont = Lato_Regular.ttf,
    BoldFont = Lato_Bold.ttf,
    ItalicFont = Lato_Italic.ttf,
    BoldItalicFont = Lato_BoldItalic.ttf
  ]
}{%
  \RequirePackage[default]{lato}
}
\renewcommand{\familydefault}{\sfdefault}

\usepackage{xcolor}
\usepackage{tikz}
\usepackage{pgfplots}
\pgfplotsset{compat=1.18}
\usetikzlibrary{arrows.meta,positioning,shapes.geometric,fit,backgrounds,calc,decorations.pathreplacing,patterns}

% Paleta institucional PLANTERR / UEFS (alinhada ao estilo planterr-slides.css / 09a)
\definecolor{etNavy}{HTML}{1F4268}
\definecolor{etNavyDark}{HTML}{152E4A}
\definecolor{etGold}{HTML}{B38A3D}
\definecolor{etGoldDark}{HTML}{8C6927}
\definecolor{etTeal}{HTML}{0E7C7B}
\definecolor{etForest}{HTML}{2D6A4F}
\definecolor{etAlert}{HTML}{9E2A2B}
\definecolor{etInk}{HTML}{26313D}
\definecolor{etMuted}{HTML}{5B6573}
\definecolor{etBorder}{HTML}{D5DEE8}

\definecolor{bgNavy}{HTML}{EDF2F7}
\definecolor{bgGold}{HTML}{F8F3E9}
\definecolor{bgTeal}{HTML}{EBF6F5}
\definecolor{bgForest}{HTML}{EEF5F0}
\definecolor{bgAlert}{HTML}{FDF2F2}
\definecolor{bgCanvas}{HTML}{FFFFFF}

\tikzset{
  >=Stealth,
  etBox/.style={
    draw=#1, fill=white, rounded corners=5pt, line width=1.15pt,
    inner sep=7pt, align=center, font=\footnotesize\sffamily, text=etInk
  },
  etTint/.style={
    draw=#1, fill=#1!8!white, rounded corners=5pt, line width=1.15pt,
    inner sep=7pt, align=center, font=\footnotesize\sffamily, text=etInk
  },
  etBadge/.style={
    fill=#1, text=white, rounded corners=9pt, inner xsep=8pt, inner ysep=3pt,
    font=\scriptsize\bfseries\sffamily
  },
  etTitle/.style={font=\normalsize\bfseries\sffamily, text=etNavy, align=center},
  etSub/.style={font=\footnotesize\sffamily, text=etMuted, align=center},
  etArrow/.style={->, >=Stealth, line width=1.35pt},
  etDArrow/.style={<->, >=Stealth, line width=1.35pt}
}


\begin{document}
\begin{tikzpicture}[font=\sffamily]

% Título
\node[etTitle, font=\large\bfseries\sffamily, text=etNavy] (title) at (7.5, 6.25)
  {Elaboração de Projetos de Investimento e Decisão Locacional};

% Linha superior: 4 etapas sequenciais
\tikzset{
  stepBox/.style={
    draw=#1, fill=white, rounded corners=6pt, line width=1.15pt,
    minimum width=3.15cm, minimum height=1.25cm, text width=2.85cm,
    align=center, font=\footnotesize\bfseries\sffamily, text=etInk
  },
  stepFocus/.style={
    draw=etGold, fill=bgGold, rounded corners=6pt, line width=1.65pt,
    minimum width=3.35cm, minimum height=1.35cm, text width=3.05cm,
    align=center, font=\footnotesize\bfseries\sffamily, text=etNavyDark
  }
}

\node[stepBox=etBorder] (s1) at (1.75, 4.95) {\textcolor{etNavy}{1. Estudo de}\\Mercado};
\node[stepBox=etBorder] (s2) at (5.45, 4.95) {\textcolor{etNavy}{2. Escala e}\\Tamanho};
\node[stepFocus]        (s3) at (9.35, 4.95) {\textcolor{etGoldDark}{3. Localização}\\Territorial};
\node[stepBox=etBorder] (s4) at (13.25, 4.95) {\textcolor{etNavy}{4. Engenharia,}\\CAPEX e OPEX};

% Badge discreto acima do foco (rounded corners <= metade da altura evita artefato nas bordas)
\node[fill=etGold, text=white, rounded corners=4pt, inner xsep=8pt, inner ysep=3pt,
  font=\scriptsize\bfseries\sffamily, anchor=south] at ([yshift=-5pt]s3.north) {DECISÃO-CHAVE};

% Setas entre etapas
\draw[etArrow, etNavy] (s1.east) -- (s2.west);
\draw[etArrow, etNavy] (s2.east) -- (s3.west);
\draw[etArrow, etGoldDark] (s3.east) -- (s4.west);

% Cards inferiores: Macrolocalização vs Microlocalização
\tikzset{
  locCard/.style={
    draw=#1, fill=white, rounded corners=6pt, line width=1.25pt,
    minimum width=6.6cm, minimum height=2.95cm, inner sep=0pt, anchor=north
  },
  locHead/.style={
    fill=#1, text=white, rounded corners=4pt,
    minimum width=6.32cm, minimum height=0.66cm,
    align=center, font=\small\bfseries\sffamily
  },
  locBody/.style={
    text width=5.95cm, align=left, font=\footnotesize\sffamily, text=etInk,
    inner sep=2pt
  }
}

\node[locCard=etNavy] (macro) at (4.0, 3.05) {};
\node[locHead=etNavy, anchor=north] (mh) at ([yshift=-0.14cm]macro.north)
  {Macrolocalização (Região / Município)};
\node[locBody, anchor=north] at ([yshift=-0.24cm]mh.south) {%
  \textbullet~Custos logísticos (insumos \textit{vs.} mercado)\\[5pt]
  \textbullet~Mão de obra e ecossistema de CT\&I\\[5pt]
  \textbullet~Incentivos regionais (SUDENE, FNE)
};

\node[locCard=etTeal] (micro) at (11.1, 3.05) {};
\node[locHead=etTeal, anchor=north] (mih) at ([yshift=-0.14cm]micro.north)
  {Microlocalização (Sítio Específico)};
\node[locBody, anchor=north] at ([yshift=-0.24cm]mih.south) {%
  \textbullet~Plano Diretor, ZEE e licenciamento\\[5pt]
  \textbullet~Infraestrutura hídrica, energética e digital\\[5pt]
  \textbullet~Externalidades socioambientais no território
};

% Setas bifurcadas da Etapa 3 para os dois cards
\coordinate (split) at (9.35, 3.68);
\draw[line width=1.35pt, etGoldDark] (s3.south) -- (split);
\draw[etArrow, etNavy] (split) -| (macro.north);
\draw[etArrow, etTeal] (split) -| (micro.north);

\end{tikzpicture}
\end{document}
